\section{Design Principles for Controlled AI-Supported Engineering}
\label{sec:design_principles}

The challenges described in Section~\ref{sec:challenges} can be translated into a set of
design principles for controlled AI-supported engineering. These principles do not define
one universal implementation architecture. Instead, they describe boundaries and control
mechanisms that proved relevant in the underlying investigation. Their potential
transferability to comparable engineering tasks is discussed in
Section~\ref{sec:discussion_transferability}.

\subsection{Start from the Engineering Task}
\label{sec:principle_task}

AI support should be designed around a sufficiently bounded engineering task. The intended
engineering result defines which information is relevant, which interpretation steps are
necessary, and which acceptance criteria can later be applied.

This task orientation is important because an AI model does not define the engineering
process in which its output will be used. The same model may support fundamentally different
activities, from information extraction to requirements analysis or model generation.
Consequently, the engineering purpose, expected output, and boundaries of AI support need to
be defined before technical implementation choices are made.

A bounded task also creates the basis for later decisions about information authority,
human review, verification, and acceptable automation.

\subsection{Establish a Controlled Information Basis}
\label{sec:principle_information_basis}

AI-supported engineering should operate on an explicitly defined information basis. The
relevant question is not only which information is available, but also which information may
justify engineering statements and which information merely supports their interpretation.

The investigated architecture therefore distinguished between \emph{Engineering Sources},
\emph{Context Sources}, and the \emph{Processing and Orchestration Context}. This distinction
made it possible to separate source-grounded evidence from information that
supported interpretation or influenced processing.

The exact source taxonomy depends on the engineering task. The general principle is that
different information roles should remain explicit and should not silently inherit the same
engineering authority. At the same time, information required to understand how a result was
produced should remain available as provenance.

\subsection{Use Probabilistic AI Where Interpretation Is Required}
\label{sec:principle_probabilistic}

Probabilistic AI is particularly useful where engineering information cannot be transformed
through explicit rules alone and instead requires semantic interpretation.

This situation occurs when source information is weakly structured, expressed in natural
language, uses inconsistent terminology, or can only be understood in relation to wider
domain context. In such cases, probabilistic processing can generate candidate
interpretations and thereby provide information for a subsequent engineering decision.

Where source and target semantics are already sufficiently formalized, however, additional
probabilistic processing may add unnecessary variability. Formal model-transformation
approaches demonstrate that explicit transformation rules can provide deterministic behavior
where the required relationships are already known \cite{Kapos.2014,Li.2022,Zhou.2024}.

The resulting principle is therefore not to maximize AI usage, but to place probabilistic
processing at those points where interpretation is actually required.

\subsection{Make Uncertainty Explicit}
\label{sec:principle_uncertainty}

Probabilistic processing should not convert uncertainty into apparently complete engineering
information without making the unresolved state visible.

In the underlying investigation, several configured perspectives were used to expose
different interpretation outcomes. Agreement could be represented as \emph{Consensus},
differences as \emph{Variance}, and jointly recognized unresolved information as
\emph{Shared Ambiguity}. These states provided additional information for subsequent human
evaluation.

The specific mechanism is not essential to the general principle. Multiple AI perspectives
represent one possible way to expose uncertainty, but other workflows may use confidence
indicators, contradiction detection, explicit unknown states, or targeted review triggers.

The important distinction is that uncertainty should remain part of the engineering state.
Agreement may strengthen the evidence available for a decision, but does not by itself prove
engineering correctness. Likewise, missing information should remain identifiable as missing
rather than being completed solely because a generative system can produce a plausible
continuation.

\subsection{Keep Engineering Authority Explicitly Assigned}
\label{sec:principle_authority}

AI-generated interpretations can support engineering decisions without themselves acquiring
engineering authority.

In the investigated architecture, the \emph{Human Engineering Authority} decided whether a
defined interpretation state could be used as authorized engineering information. A separate
\emph{Model Authority} addressed the structural placement and representation of already
authorized content in the target model.

This separation emerged because engineering meaning and model representation constituted
different decision objects. A statement may be accepted as technically meaningful without
this automatically determining how it should be represented in a model or downstream
artifact.

The exact authority roles are context dependent. Some engineering tasks may require fewer
decision boundaries, while others may require additional domain, quality, or release
authorities. The transferable principle is that consequential state transitions should have
an explicitly assigned decision authority instead of allowing probabilistic processing to
implicitly determine their outcome.

Authorization should refer to a clearly identifiable information state. If that state is
materially changed, the previous authorization should not automatically be assumed to remain
valid.

\subsection{Constrain Downstream Processing Once Decisions Are Resolved}
\label{sec:principle_constraints}

Once the relevant engineering meaning, decision authority, and representation rules are
sufficiently resolved, subsequent processing should become increasingly rule-bound and
reproducible.

The investigated architecture used this separation to avoid repeated probabilistic
reinterpretation of already authorized information. Authorized engineering content was
carried forward into a controlled engineering state and processed through explicit modeling
and transformation rules before the final artifact was produced.

This principle reduces the number of points at which probabilistic behavior can change an
already established engineering decision. It also improves the ability to reproduce and
verify downstream processing.

A related consequence is fail-closed behavior. If required evidence, authorization, or a
necessary transformation rule is missing, processing should expose the unresolved condition
rather than silently infer a continuation. This does not imply that every workflow must stop
at the same type of event. It means that unresolved engineering decisions should not be
hidden by automatic completion.

\subsection{Preserve Provenance and Traceability}
\label{sec:principle_traceability}

Controlled AI-supported engineering requires more than storing the final result. The
relationships between source information, AI-generated interpretations, human decisions,
intermediate engineering states, and resulting artifacts should remain reconstructable.

Provenance describes where information originates and how it was derived
\cite{MoreauMissier2013,Simmhan2005}. For AI-derived interpretations, this also includes
the relevant processing conditions under which they were produced. Traceability connects
the relevant engineering states and decision objects throughout the workflow. Together, they make it possible to determine
why a result exists, which information contributed to it, and which decisions authorized its
use.

This becomes particularly important in multi-source scenarios. Information from several
Engineering Sources may contribute to one resulting engineering state, but integration
should not remove the ability to identify the individual source contribution and its
associated decision history.

Persisted intermediate states also support change analysis. If a source changes, affected
interpretations, decisions, and downstream artifacts can be identified rather than treating
the final result as an isolated output.

\subsection{Assess Results Separately from Their Generation}
\label{sec:principle_verification}

Generation and the subsequent assessment of the resulting engineering artifact should
remain distinct functions.

An AI system may produce a syntactically plausible or formally valid artifact while the
artifact still contains engineering-relevant errors, incomplete relationships, or unsupported
content \cite{Topcu2025,Cibrian.2025,Stein2026}. The fact that the generation process
completed successfully is therefore not sufficient evidence that the result satisfies its
engineering purpose.

The required assessment mechanisms should instead be selected according to the target
artifact and the engineering task. Depending on the application, this may include syntax
checks, structural validation, semantic constraints, consistency checks, model-tool
validation, human model review, or further domain-specific acceptance criteria.

The investigated architecture treated technical validation, human model review, and final release
as distinct concerns. This separation prevents successful generation from being interpreted
as automatic approval and creates an explicit assurance boundary before the resulting
artifact is used in subsequent engineering work.

Taken together, these principles describe a controlled division of responsibilities:
probabilistic AI supports interpretation where ambiguity exists; engineering authority
governs consequential decisions; rule-bound processing handles sufficiently determined
states; provenance and traceability preserve the processing history; and separate
assessment provides evidence about the resulting artifact.
